\subsection{Листинг численного метода}

Для реализации блок-схем воспользуемся языком программирования C++.

\begin{lstlisting}[caption=Вычисление значения интерполирующей функции \( F(x) \) заданной точности.]
double interpolate_by_lagrange(int f, double a, double b, double x) {
  function<double(double)> func = get_function(f);
  int n = 1;
  std::vector<double> roots(n);
  fill_chebyshev_polynomial_roots(n, a, b, roots.data());
  double lagrange_old, lagrange_current = generate_lagrange_polynomial(n, func, roots.data())(x);
  do {
    n++;
    lagrange_old = lagrange_current;
    roots.reserve(n);
    fill_chebyshev_polynomial_roots(n, a, b, roots.data());
    lagrange_current = generate_lagrange_polynomial(n, func, roots.data())(x);
  } while (abs(lagrange_current - lagrange_old) > 1e-4 && n < 100);
  return lagrange_current;
}
\end{lstlisting}

\begin{lstlisting}[caption=Построение интерполяционного многочлена в форме Лагранжа \( L_n(t) \).]
function<double(double)> generate_lagrange_polynomial(const int n, const function<double(double)> func, const double grid[]) {
  return [n, func, grid](double x) {
    double result = 0;
    for (int i = 0; i < n; i++) {
      double current_term = func(grid[i]);
      for (int j = 0; j < n; j++) {
        if (j != i) {
          current_term *= x - grid[j];
          current_term /= grid[i] - grid[j];
        }
      }
      result += current_term;
    }
    return result;
  };
}
\end{lstlisting}

\begin{lstlisting}[caption=Поиск корней многочлена Чебышева первого рода \( T_n(t) \).]
void fill_chebyshev_polynomial_roots(const int n, const double a, const double b, double roots[]) {
  for (int k = 1; k <= n; k++) {
    roots[k - 1] = 0.5 * (a + b + (b - a) * cos(M_PI * (2 * k - 1) / (2.0 * n)));
  }
}
\end{lstlisting}